%=============================================================================
% Interpretable Statistical Modeling for Real-Time Sepsis-3 Onset Prediction
% in MIMIC-IV: A Dynamic Landmarking Approach with Explicit Assumption Testing
% and Informative-Missingness Analysis
%
% IEEE conference format.  Compile:
%   pdflatex sepsis_realtime_ieee
%   bibtex   sepsis_realtime_ieee
%   pdflatex sepsis_realtime_ieee
%   pdflatex sepsis_realtime_ieee
% Bibliography file: sepsis_realtime_refs.bib
%=============================================================================
\documentclass[conference]{IEEEtran}
\IEEEoverridecommandlockouts

\usepackage{amsmath,amssymb,amsfonts}
\usepackage{graphicx}
\usepackage{booktabs}
\usepackage{array}
\usepackage{textcomp}
\usepackage{xcolor}
\usepackage[hidelinks]{hyperref}
\usepackage{csquotes}
\usepackage[style=ieee]{biblatex}
\addbibresource{sepsis_realtime_refs.bib}

\begin{document}

\title{Interpretable Statistical Modeling for Real-Time Sepsis-3 Onset
Prediction in MIMIC-IV: A Dynamic Landmarking Approach with Explicit
Assumption Testing and Informative-Missingness Analysis}

\author{\IEEEauthorblockN{Abdul Ahad}
\IEEEauthorblockA{\textit{Department of Computer Science \& Applied Physics} \\
\textit{Atlantic Technological University (ATU)}\\
Galway, Ireland \\
ahadarif.1998@gmail.com}
}

\maketitle

%------------------------------------------------------------------ ABSTRACT
\begin{abstract}
Sepsis is a leading cause of intensive-care mortality, and early recognition
improves survival. Recent prediction research is dominated by machine learning
and, on the MIMIC-IV database specifically, by \emph{post-diagnosis}
prognostic modeling (mortality nomograms) rather than \emph{pre-onset} early
detection. The one statistical real-time onset benchmark, TREWScore, was built
on the older MIMIC-II database in 2015 and has not been revisited on MIMIC-IV
with purely statistical, interpretable tools. We develop and internally
validate a fully interpretable statistical pipeline for real-time Sepsis-3
\emph{onset} prediction on MIMIC-IV (65{,}241 adult ICU stays), reaching
continuously updated risk through a dynamic \emph{landmark supermodel} rather
than any neural or ensemble method. Three methodological commitments
distinguish the work. First, the proportional-hazards assumption underlying
Cox regression is \emph{formally tested} (Schoenfeld residuals) and, because it
fails globally, a time-varying reformulation is fitted in response, a step
routinely omitted in the MIMIC-IV Cox literature. Second, incident (new-onset)
sepsis is isolated from prevalent sepsis using an hour-6 landmark, exposing
that standard bedside scores are non-discriminative to \emph{inverse} for this
task (SOFA AUROC $0.36$, qSOFA $0.49$, SIRS $0.59$), whereas a multivariable
landmark model attains AUROC $0.775$. Third, an informative-missingness
analysis shows that \emph{whether a laboratory test was ordered} is a stronger
predictor than its value, quantifying a label-recognition ceiling intrinsic to
retrospectively labeled ICU data. The pipeline is closed-form, transportable,
and pre-wired for external validation on eICU-CRD, directly targeting the
internal/external validation asymmetry that recent systematic reviews identify
as the field's central unresolved issue. All extraction, modeling, and
evaluation code is released for reproducibility.
\end{abstract}

\begin{IEEEkeywords}
sepsis prediction, survival analysis, landmarking, proportional hazards,
informative missingness, MIMIC-IV, interpretable models, electronic health
records
\end{IEEEkeywords}

%------------------------------------------------------------------ INTRO
\section{Introduction}
Sepsis, defined by the Sepsis-3 consensus as life-threatening organ
dysfunction caused by a dysregulated host response to infection
\cite{singer2016sepsis3}, remains a major driver of ICU mortality worldwide.
Because early treatment materially improves outcomes, predicting onset ahead of
clinical recognition is of direct clinical value. Two modeling traditions
address this problem. The first, and by far the more actively published, is
machine learning; the second, older and still foundational, is classical
statistics, namely Cox proportional-hazards regression, logistic regression, and
trajectory modeling. Statistical models retain a specific advantage for
real-time clinical use: they yield directly interpretable hazard and odds
ratios, they can be checked against formal assumptions, and they can be
benchmarked transparently against established rule-based scores (SIRS, qSOFA,
NEWS).

On the MIMIC-IV database \cite{johnson2023mimiciv}, however, the statistical
literature has developed unevenly. A companion review of eighteen peer-reviewed
statistical studies found the body of work methodologically mature but
concentrated on the wrong half of the problem: most studies predict mortality
\emph{after} sepsis is diagnosed rather than onset itself, the
proportional-hazards assumption is rarely tested before a Cox model is fitted,
and external validation is largely absent from the onset-prediction literature
specifically. The single statistical real-time onset benchmark, the TREWScore
of \textcite{henry2015trewscore}, was fitted on the older
MIMIC-II database in 2015; \textcite{liaw2019mortality} used an
independent prospective cohort. Neither used MIMIC-IV.

This paper develops a statistical pipeline that addresses these limitations
within a single framework and, critically, states its \emph{research gap},
\emph{research question}, and \emph{contributions} explicitly (Sections
\ref{sec:gap}--\ref{sec:contrib}) so that the novelty relative to prior work is
auditable. We deliberately restrict the modeling to classical statistics, both
for interpretability and to establish a rigorous, transportable baseline
against which machine-learning approaches can later be compared.

%------------------------------------------------------------------ RELATED
\section{Related Work}
\label{sec:related}

\subsection{Statistical onset prediction}
TREWScore \cite{henry2015trewscore} fitted a Cox model to routinely collected
MIMIC-II vitals and labs to anticipate septic shock, reaching AUC $0.83$ with a
median 28-hour lead time, the foundational statistical benchmark for real-time
sepsis prediction. \textcite{liaw2019mortality} proposed a
longitudinal, time-varying Cox extension tracking evolving mortality risk from
repeatedly measured indicators. \textcite{kasal2004cox}
compared a standard Cox model against Gray's flexible survival model in severe
sepsis and found that hazards are frequently \emph{non-proportional}, a
methodological caution largely ignored by later MIMIC-IV work.

\subsection{Prognostic modeling on MIMIC-IV}
A recurring MIMIC-IV pattern couples LASSO variable selection with
multivariable logistic regression, presented as a nomogram, to predict
mortality after diagnosis; several such studies include external validation on
eICU-CRD \cite{yuan2023nomogram}. Trajectory approaches such as group-based
trajectory modeling (GBTM) of SOFA scores \cite{yang2022gbtm} model the
evolving course but again target mortality, not onset. This sub-literature is
well validated but prognostic, not predictive of onset.

\subsection{Benchmarks, validation, and reviews}
The PhysioNet/CinC 2019 Challenge \cite{reyna2020physionet} established a shared
early-prediction benchmark and a clinical-utility metric.
\textcite{fleuren2020ml} reviewed machine-learning sepsis prediction and reported
large between-study heterogeneity. Most pointedly,
\textcite{wang2025srpm} reviewed 91 sepsis real-time prediction models and found
that median AUROC falls from $0.886$ internally to $0.783$ under full external
validation, evidence that the internal/external gap, not raw internal
discrimination, is the field's central weakness.

\subsection{Landmarking and informative missingness}
Dynamic prediction by landmarking \cite{vanhouwelingen2007landmark} and its
supermodel formulation \cite{fries2023dynamiclm} pool multiple prediction
windows to emit continuously updated risk from a single interpretable model,
providing a statistical route to real-time prediction without recurrent
networks. Separately, informative missingness (the fact that \emph{whether} a
test is ordered encodes clinician suspicion) is well documented in EHR
modeling \cite{groenwold2020missingness,singh2021missingness}. Both threads are
central to our design and, to our knowledge, have not been combined for
statistical Sepsis-3 onset prediction on MIMIC-IV.

%------------------------------------------------------------------ GAP
\section{Research Gap}
\label{sec:gap}
Synthesizing Section \ref{sec:related}, three gaps are specific and actionable:

\begin{itemize}
\item \textbf{G1 (Onset vs.\ prognosis).} MIMIC-IV statistical work overwhelmingly
predicts mortality \emph{after} diagnosis. Genuine pre-onset, real-time
detection with purely statistical tools has not been revisited on MIMIC-IV
since TREWScore's 2015 MIMIC-II study.
\item \textbf{G2 (Untested proportional hazards).}
\textcite{kasal2004cox} showed sepsis hazards are frequently non-proportional, yet
no reviewed MIMIC-IV Cox study reports testing the assumption, weakening
confidence in the reported hazard ratios.
\item \textbf{G3 (Validation asymmetry).} External validation is common in the
prognostic nomogram literature but essentially absent from onset-prediction and
trajectory studies, precisely where reviews show performance degrades most
\cite{wang2025srpm}.
\end{itemize}

%------------------------------------------------------------------ RQ
\section{Research Question and Objectives}
\label{sec:rq}
These gaps motivate one primary research question and three sub-questions.

\vspace{2pt}
\noindent\textbf{RQ.} \emph{Can a purely statistical (Cox proportional-hazards
/ time-varying / landmark) model, using routinely available real-time vitals
and labs in MIMIC-IV, predict impending Sepsis-3 onset ahead of clinical
recognition, with performance comparable to or exceeding the rule-based scores
SIRS, qSOFA, and NEWS?}

\begin{itemize}
\item \textbf{SQ1.} Which routinely measured variables are the strongest
statistically significant predictors of onset, and \emph{does the
proportional-hazards assumption hold} for the fitted Cox model?
\item \textbf{SQ2.} Does modeling trajectories through a time-varying / landmark
extension improve early detection over a static single-time-point model?
\item \textbf{SQ3.} How does discriminative performance compare with SIRS,
qSOFA, and NEWS, and does it survive the transition from internal to external
(eICU-CRD) evaluation?
\end{itemize}

\noindent The null hypothesis for SQ1 is that proportional hazards holds for at
least half of the included covariates; its rejection promotes the time-varying
reformulation (Objective under G2) from a secondary to a primary contribution.

%------------------------------------------------------------------ CONTRIB
\section{Contributions}
\label{sec:contrib}
This work makes the following contributions, each mapped to the gap it closes.

\begin{enumerate}
\item[\textbf{C1.}] \textbf{First interpretable statistical real-time Sepsis-3
\emph{onset} model on MIMIC-IV (closes G1).} We reach continuously updated,
hourly risk through a dynamic landmark supermodel
\cite{vanhouwelingen2007landmark,fries2023dynamiclm}, a classical-statistics
alternative to recurrent or boosted models, and explicitly separate
\emph{incident} from \emph{prevalent} sepsis using an hour-6 landmark, since
$\sim$34\% of septic stays are already septic at admission and would otherwise
inflate apparent performance.

\item[\textbf{C2.}] \textbf{Explicit proportional-hazards assumption testing with
a triggered time-varying reformulation (closes G2).} We report the Schoenfeld
residual test for the sepsis-onset Cox model; the global test \emph{fails}, and
a time-varying Cox model is fitted in direct response. This diagnostic step is
routinely omitted in the MIMIC-IV Cox literature, so reporting it, and acting
on it, is itself a methodological contribution.

\item[\textbf{C3.}] \textbf{Reframing the task: bedside severity scores do not
predict incident onset.} On the correct incident-onset target, SOFA is
\emph{inversely} discriminative (AUROC $0.36$), qSOFA is near chance ($0.49$),
and SIRS is weak ($0.59$), while a multivariable landmark model reaches
$0.775$. This quantifies a distinction the field routinely blurs: severity
scores detect \emph{established} organ failure, not impending onset.

\item[\textbf{C4.}] \textbf{Informative-missingness characterization and a
label-recognition ceiling.} We show that binary ``was this lab ordered''
indicators are stronger predictors than the imputed values themselves,
connecting the result to the ``curse of knowing''
\cite{groenwold2020missingness,singh2021missingness} and quantifying an upper
bound on what \emph{any} model can achieve when the onset label is itself gated
by clinician recognition. This is an honest, reusable characterization result,
not a leaderboard claim.

\item[\textbf{C5.}] \textbf{A reproducible, transportable pipeline pre-wired for
external validation (targets G3).} Because the nomogram and supermodel are
closed-form, their coefficients transport directly to a second database; the
pipeline is structured so that eICU-CRD \cite{pollard2018eicu} validation
requires only re-running the cohort/onset/landmark stages, directly addressing
the validation asymmetry that reviews flag as decisive \cite{wang2025srpm}.
\end{enumerate}

%------------------------------------------------------------------ METHODS
\section{Data and Methods}
\label{sec:methods}

\subsection{Cohort and labeling}
From MIMIC-IV v3.1 we build an adult ICU cohort (age $\ge 18$, first ICU stay
only, length of stay $\ge 4$\,h), yielding 65{,}241 stays. Suspected infection
follows \textcite{seymour2016criteria}: antibiotic orders
paired with cultures within the Sepsis-3 windows, with surveillance swabs (e.g.
MRSA screens) excluded to avoid inflating the suspicion rate. Sepsis-3 is
suspected infection together with an acute SOFA rise $\ge 2$; prevalence is
37.1\%, with septic shock 9.2\%.

\subsection{Hourly panel and onset labeling}
Every stay is resampled to an hourly grid ($\sim$3.1\,M stay-hours). For each
variable we retain the raw value, a forward-filled value (what the clinician
currently knows), and a time-since-last-measured counter. Hourly SOFA locates
the exact onset hour and defines a per-hour ``onset within 12\,h'' label. An
hour-6 landmark yields 48{,}829 at-risk stays, of which 3{,}330 develop onset
within 12\,h.

\subsection{Models}
Baselines are the rule-based scores SIRS, qSOFA, and NEWS computed directly
from vitals. Statistical models comprise: (i) a static Cox model with Schoenfeld
testing and a time-varying reformulation on failure; (ii) an elastic-net
penalized logistic \emph{nomogram} at the hour-6 landmark using an
informative-missingness design (median impute each lab, then add a ``was
measured'' indicator); (iii) a \emph{landmark supermodel} stacked across
landmarks at 6, 12, 18, 24, 36, and 48\,h, so one model emits hourly-updating
risk; (iv) GEE with cluster-robust standard errors for within-patient
correlation; (v) GBTM of first-24\,h SOFA trajectories linked to onset and
mortality; and (vi) confounder-adjusted causal estimates (PSM, IPW, augmented
IPW, doubly robust) for early antibiotics. All randomness uses a fixed seed for
reproducibility.

\subsection{Evaluation}
Discrimination uses time-dependent AUROC; calibration uses Brier score and
calibration curves; clinical utility uses the PhysioNet/CinC 2019 utility score
\cite{reyna2020physionet} and an alarm-burden curve (false alarms per
patient-day at fixed sensitivity). Subgroup calibration across sex and age
serves as an internal fairness check and as a stand-in until eICU-CRD external
validation is executed.

%------------------------------------------------------------------ RESULTS
\section{Results}
\label{sec:results}

\subsection{Benchmarking on the incident-onset task}
Table \ref{tab:bench} reports discrimination at the hour-6 landmark for
incident onset within 12\,h. The rule-based scores are non-discriminative to
inverse; the statistical landmark models clearly exceed them.

\begin{table}[t]
\caption{Discrimination at the hour-6 landmark (incident onset within 12\,h).}
\label{tab:bench}
\centering
\begin{tabular}{lc}
\toprule
\textbf{Model / Score} & \textbf{AUROC} \\
\midrule
SOFA at hour 6                              & 0.36 (inverse) \\
qSOFA at hour 6                             & 0.49 \\
SIRS at hour 6                              & 0.59 \\
Multivariable logistic + trajectory slopes & 0.755 \\
Full-feature elastic-net nomogram (C1)     & \textbf{0.775} \\
\bottomrule
\end{tabular}
\end{table}

The trajectory-slope terms add statistically significant signal over the
static model (DeLong $p \approx 3\times10^{-7}$), supporting SQ2.

\subsection{Assumption testing (SQ1)}
The global Schoenfeld test rejects proportional hazards for the static onset
Cox model, so the null of SQ1 is rejected and a time-varying Cox model is
fitted in response, satisfying C2. This confirms, on MIMIC-IV, the
non-proportionality that \textcite{kasal2004cox} reported
elsewhere.

\subsection{Real-time behavior and honest ceilings}
Table \ref{tab:analyses} summarizes the deployable model and its supporting
analyses. Two results are deliberately framed as \emph{characterizations}
rather than wins: the strongest nomogram predictors are missingness indicators
(C4), and the PhysioNet utility is modest because the onset label is gated by
clinician recognition, an intrinsic ceiling rather than a model deficiency.

\begin{table}[t]
\caption{Real-time model and supporting analyses (internal MIMIC-IV).}
\label{tab:analyses}
\centering
\begin{tabular}{p{0.30\columnwidth}p{0.58\columnwidth}}
\toprule
\textbf{Component} & \textbf{Headline result} \\
\midrule
Nomogram        & AUROC 0.775; ``was-measured'' indicators dominate (C4) \\
Supermodel      & Hourly-updating risk; AUROC(s) 0.65--0.77 (C1) \\
GEE / TV-Cox    & Within-patient corr.\ 0.68; cluster-robust SE up to $4.5\times$ naive \\
Utility         & Normalized utility 0.29; 90\% sensitivity $\Rightarrow$ 6.5 false alarms/patient-day \\
Subgroups       & Calibration ratio 0.94--1.06 across sex/age \\
GBTM            & 3 latent SOFA-trajectory classes; steepest-worsening $\sim 5\times$ mortality HR \\
Causal          & Early antibiotics ($\le 3$\,h) associated with lower mortality across PSM/IPW/AIPW/DR \\
\bottomrule
\end{tabular}
\end{table}

%------------------------------------------------------------------ DISCUSSION
\section{Discussion}
\label{sec:discussion}
The central empirical finding is that the ceiling on incident-onset prediction
is set by the label's dependence on clinician recognition (when cultures and
antibiotics are ordered) rather than by model class. This is why the strongest
predictors are missingness indicators and why a purely statistical model is
competitive with far more complex approaches: no neural network is required to
approach a ceiling imposed by the labeling process itself. The result aligns
with, and sharpens, the informative-missingness literature
\cite{groenwold2020missingness,singh2021missingness} by placing it at the
center of a real-time onset task rather than treating it as a nuisance.

Methodologically, reporting and acting on the proportional-hazards test (C2)
restores a diagnostic step the MIMIC-IV Cox literature skips. Reframing the
benchmark (C3) corrects a common conflation of severity scoring with onset
prediction: SOFA's inverse discrimination is not an anomaly but a direct
consequence of scoring \emph{established} dysfunction.

\subsection{Threats to validity}
The onset label depends on documentation timing; the assumed-zero SOFA baseline
may slightly overcount organ dysfunction; qSOFA uses mean arterial pressure as a
proxy for systolic pressure; and the time-varying Cox and HMM analyses use
sampled subsets for tractability. Most importantly, all results are
\emph{internal} to MIMIC-IV. Given that reviews report a substantial
internal-to-external AUROC drop \cite{wang2025srpm}, the external validation in
C5 is required before any clinical-generalizability claim; it is designed but
not yet executed, and is the primary item of future work.

%------------------------------------------------------------------ CONCLUSION
\section{Conclusion and Future Work}
\label{sec:conclusion}
We presented an interpretable, purely statistical pipeline for real-time
Sepsis-3 onset prediction on MIMIC-IV that (C1) reaches hourly-updating risk via
a landmark supermodel, (C2) tests and acts on the proportional-hazards
assumption, (C3) shows bedside scores fail on the incident-onset task, (C4)
characterizes an informative-missingness ceiling, and (C5) is transportable and
pre-wired for external validation. The immediate next step is eICU-CRD external
validation \cite{pollard2018eicu} to quantify the internal/external gap for this
statistical model; subsequent work will use the same landmark substrate to
benchmark machine-learning approaches against this interpretable baseline on
identical splits.

%------------------------------------------------------------------ BIB
\printbibliography

\end{document}
